\documentclass{article}


% if you need to pass options to natbib, use, e.g.:
%     \PassOptionsToPackage{numbers, compress}{natbib}
% before loading neurips_2024

% to compile a preprint version, e.g., for submission to arXiv, use the
% [preprint] option; for the anonymized submission version, remove it.
\usepackage[preprint]{neurips_2024}

% to compile a camera-ready version, add the [final] option, e.g.:
%     \usepackage[final]{neurips_2024}


\usepackage[utf8]{inputenc} % allow utf-8 input
\usepackage[T1]{fontenc}    % use 8-bit T1 fonts
\usepackage{hyperref}       % hyperlinks
\usepackage{url}            % simple URL typesetting
\usepackage{booktabs}       % professional-quality tables
\usepackage{amsfonts}       % blackboard math symbols
\usepackage{amsmath}        % math environments
\usepackage{nicefrac}       % compact symbols for 1/2, etc.
\usepackage{microtype}      % microtypography
\usepackage{xcolor}         % colors
\usepackage{graphicx}       % include figures


\title{Kepler-Encoder-V1 Technical Report:\\
       Towards a Multimodal Embedding Model for Robots}


% The \author macro works with any number of authors. Use \And to let LaTeX
% break lines, or \AND to force a line break at that point.
\author{%
  Author Name\thanks{Correspondence: \texttt{email@example.com}} \\
  Affiliation \\
  \texttt{email@example.com} \\
  % \And
  % Coauthor \\
  % Affiliation \\
  % \texttt{email@example.com} \\
}


\begin{document}


\maketitle


\begin{abstract}
  We present Kepler-Encoder-V1, a multimodal embedding model for robotics that
  maps heterogeneous robot observations and actions into a shared latent space.
  The model couples a Perceiver~IO backbone~\citep{perceiverio2021} with a
  joint-embedding predictive training objective~\citep{lejepa2025}, and is
  trained on large-scale real-robot interaction data~\citep{rh20t2023}.
  % TODO: summarize the setup, key quantitative results, and headline finding.
  This report describes the architecture, training objective, data pipeline,
  and preliminary evaluation of the V1 model.
\end{abstract}


\section{Introduction}
\label{sec:intro}
% Motivation: why a single multimodal embedding space for robots.
Robots perceive and act through many heterogeneous channels---RGB video,
proprioception, force/torque, and language---and learning a single
representation that unifies them remains an open problem. Recent work has
explored multimodal representations for prediction and
control~\citep{multimodal2019} and self-supervised predictive pre-training for
motor control~\citep{spatiotemporal2024}. Kepler-Encoder-V1 targets a shared
embedding space usable across these modalities and across robot embodiments.
% TODO: state the concrete problem, contributions, and paper roadmap.

\paragraph{Contributions.}
% TODO: enumerate contributions.
\begin{itemize}
  \item A multimodal encoder architecture based on Perceiver~IO
        \citep{perceiverio2021}.
  \item A training recipe built on the LeJEPA joint-embedding objective
        \citep{lejepa2025}.
  \item An evaluation on real-robot data~\citep{rh20t2023}. % TODO: results.
\end{itemize}


\section{Related Work}
\label{sec:related}

\paragraph{Joint-embedding and predictive representations.}
Joint-embedding predictive architectures learn representations by predicting
latent targets rather than reconstructing pixels. This paradigm has been
applied to functional brain signals~\citep{jmaefmri2024}, multivariate time
series~\citep{mtsjepa2026}, and robotic motor
control~\citep{spatiotemporal2024}. Our training objective follows
LeJEPA~\citep{lejepa2025}.

\paragraph{Robot representation and action tokenization.}
A growing line of work tokenizes robot actions for use with sequence models,
including VQ-based action tokenizers~\citep{vqvla2025}, frequency-space action
tokenization~\citep{fast2025,stablefast2025}, learned action
codecs~\citep{actioncodec2026}, and heterogeneous pre-trained
transformers~\citep{hpt2024}. Motion instruction with inertial tokens has also
been explored~\citep{mojito2025}.

\paragraph{Image and video autoencoders.}
Compression-oriented and diffusion-based video autoencoders inform our
tokenization of the visual stream~\citep{adaptive1dvdae2026,rvmae2026,
videoae2601,videoae2604,multimodalcompression2022}. Neural codecs have
additionally been repurposed as tokenizers for adjacent signal
domains~\citep{biosignaltokenizers2025}.

\paragraph{Robot interaction datasets.}
We build on large-scale real-robot datasets~\citep{rh20t2023}, and situate our
data choices relative to humanoid~\citep{humanoidseveryday2025}, large-scale
manipulation~\citep{agibotworld2025}, mobile
manipulation~\citep{airoamoma2025}, and contact-rich force-aware
collections~\citep{forcevla22026}.


\section{Method}
\label{sec:method}

\subsection{Architecture}
\label{sec:arch}
% TODO: describe the Perceiver IO encoder, latent array, cross-attention, and
% per-modality input/output adapters.
Kepler-Encoder-V1 uses a Perceiver~IO backbone~\citep{perceiverio2021} that
cross-attends heterogeneous input arrays into a fixed-size latent array,
enabling a modality-agnostic encoder.

\subsection{Training Objective}
\label{sec:objective}
% TODO: state the LeJEPA objective, target/context construction, and loss.
Training follows the LeJEPA joint-embedding predictive
objective~\citep{lejepa2025}.

\subsection{Multimodal Tokenization}
\label{sec:tokenization}
% TODO: describe per-modality tokenizers (video, proprioception, action, force).


\section{Data}
\label{sec:data}
% TODO: describe the RH20T-based training corpus, preprocessing, splits.
The primary training corpus is RH20T~\citep{rh20t2023}.
% TODO: table of dataset statistics.


\section{Experiments}
\label{sec:experiments}
% TODO: describe evaluation protocol, baselines, and metrics.

\subsection{Setup}
\label{sec:setup}
% TODO.

\subsection{Results}
\label{sec:results}
% TODO: main results table and analysis.


\section{Discussion and Limitations}
\label{sec:discussion}
% TODO.


\section{Conclusion}
\label{sec:conclusion}
% TODO.


%%%%%%%%%%%%%%%%%%%%%%%%%%%%%%%%%%%%%%%%%%%%%%%%%%%%%%%%%%%%%%%%%%%%%%%%%%%%%%%%%%
% References
%%%%%%%%%%%%%%%%%%%%%%%%%%%%%%%%%%%%%%%%%%%%%%%%%%%%%%%%%%%%%%%%%%%%%%%%%%%%%%%%%%
\bibliographystyle{plainnat}
\bibliography{references}


\end{document}
